% Modernised reading — fonds Alexandre Grothendieck, shelfmark 156-3, pages 1-40.
%
% This is an INTERPRETATION, not a transcription. Everything the critical
% apparatus carried moves into footnotes. In any dispute of substance the
% transcription governs: whoever cites Grothendieck must cite batch-01.fr.tex
% (pages 1-20) and batch-02.fr.tex (pages 21-40).
%
% Pass: Opus 5.5 (claude-opus-5-5), 2026-10-10 — first pass, unchecked against
% the pages by a human. The folder was read whole: two batches, forty pages,
% all mathematics in his hand, dated 8 June 1986 (page 1) and 9 June (page 21).
%
% What the folder is. Chapter III (« GF III ») of Vers une géométrie des
% formes, « Réseaux via découpages ». It proposes a third axiomatics of the
% « réseaux » that chapter I (156-1) defined by segments and chapter II
% (156-2) announced without reaching: a set of lieux, the pairs of lieux in
% « position relative modérée » (0-sphères), and for each finite moderate
% part the decomposition of its moderate complement into components. The
% axiomatics is tested only in dimension one, on the « tronçon ordonné », and
% the folder becomes a study of the tame (« modérées ») parts of a tronçon and
% of the algebra they form. It ends (page 40) by restating the dimension-one
% problem and leaving it, as presumably classical (« Hilbert ? »).
%
% Notation fixed for the whole document. 𝓛 the set of lieux, and from page 6
% on the tronçon (the page writes a plain L from page 21; kept 𝓛 here).
% 𝔖_0 ⊂ 𝔓_2(𝓛) the 0-sphères; Modf(𝓛) the finite moderate parts; C(T) the
% moderate complement; R_T its equivalence; Ȧ (dot) is used ONLY for the
% abstract frontier datum (c) of page 3. Drap_2, Drap the comparable pairs and
% finite chains. 𝓡(a,b,c) strict betweenness, 𝓡̄ non-strict. ∂T the absolute
% boundary of page 25 (extreme points of T); Fr X the moderate frontier of a
% moderate part (the page's Ẋ on pages 18-20, its Fr from page 29, and the
% abstract Fr of page 33); Fr_top the frontier for the canonical topology (the
% page's Ṫ on page 25). ≺, ≪, Tr(𝓛), X_Φ, π_0, X', 𝓜 = Mod(𝓛), 𝓜_f, X̄
% (moderate closure), X°, 𝓜_A are his.
%
% Substantive departures from the pages, each footnoted where it occurs:
% — page 1: the bracket « En termes de l'axiomatique par les segments » is read
%   only in part; the reading gives its sense and says so.
% — page 2: « s ≠ t » in the definition of a moderate pair is read s = t.
% — page 5: « U ⊂ V i.e. T ⊂ S » — only T ⊂ S ⇒ C(S) ⊂ C(T) is used.
% — page 10: the inner components are the open ]s_i, s_{i+1}[ (the page has
%   brackets that read closed).
% — page 11: the margin's « or {a,b,c} of cardinal ≤ 2 » is replaced by the
%   exact relation 𝓡̄(a,b,c) ⇔ 𝓡(a,b,c) or b ∈ {a,c}.
% — page 13: « trois » éléments is four.
% — page 16: the claim that the largest class 𝒯 is the class of tronçons is
%   verified (footnote).
% — page 18: uniqueness of the decomposition is proved (footnote): the
%   components are the maximal pseudo-tronçons contained in X.
% — page 22: T_0 is stated by the page for the lieux of a faille; for a
%   general tronçon (his open question) it fails: ℝ with a doubled origin.
% — page 23: the topology generated by both families is the least fine of
%   those finer than both; the list of opens is a sub-basis; « [a,b] » read ∁[a,b].
% — page 25: ∂𝓛_{≥a} = {a} only when 𝓛 has no greatest element; the
%   disjointness in ∂T = (Fr_top T ∩ T) ⊔ (T ∩ ∂𝓛) fails when T is a point of ∂𝓛.
% — page 26: « [b,a[ » read [a,b[.
% — page 27: T ≺ T' is taken with the moderate-position and disjointness
%   hypotheses of the margin; « T a un plus grand élément » is read « T a une
%   borne supérieure a = Sup T, extrémité de son adhérence ».
% — page 28: the interstitial S is the pseudo-tronçon with ends a, b, not the
%   set of x with T ⊂ 𝓛_{<x}, 𝓛_{>x} ⊃ T'; the two differ for a doubled
%   point. Same for Cor 1. « 𝓛_{≤x} » of page 27 read 𝓛_{<x}.
% — page 31: « minoré » / « majoré » read strictly; the count n' fails for n = 0.
% — page 33: « Fr X ∪ Fr X » read Fr X ∪ Fr Y; X ∨ Y = X + Y + XY (the page's
%   minus sign is the same in characteristic 2).
% — page 34: 1_{𝓜_A} = A (page: X); 𝓜_Y read 𝓜_A.
% — page 38: « ∂T_i = ∅ ⇒ T_i = 𝓛 » holds for the moderate frontier, not for
%   the absolute boundary when 𝓛 is a segment.
% — page 39: ]t_1,t_3[ read ]t_2,t_3[; the atoms cover 1 in 𝓜, not 𝓛 as a set.
%
% Modern names given and footnoted as ours: relation d'intermédiarité
% (betweenness: Pasch 1882, Hilbert 1899, Huntington–Kline 1917; for partial
% orders Altwegg 1950), topologie de l'ordre, axiome T_0 and espace sobre (the
% latter his own term, SGA 4), algèbre d'intervalles (Mostowski–Tarski 1939),
% topologie modérée and structures o-minimales (van den Dries 1984,
% Pillay–Steinhorn 1986; van den Dries 1998), algèbre de Boole partielle
% (Kochen–Specker 1965-1967).
%
% What the folder announces and never establishes: the axioms themselves —
% the data (a), (b), (c) of pages 1-3 are given, the axioms that should make
% (c) follow from (b) and (b) from its 0-sphere part are « prévus » and never
% written beyond Dec_1, Dec_2 and the axiom of page 4; the conditions on
% 𝓡(a,b,c), and on (𝔖_0, R_ε), characterising a tronçon (page 16 (2), page
% 40); the end of the locally closed theorem of page 20 (overwritten);
% properties 1)-8) of pages 32-33, stated without proof; soberness (page 22).
\documentclass[11pt,a4paper]{article}
\input{../preamble/grothendieck}

\folder{156-3}
\pages{1}{40}
\dating{1986}
\watermark{Édition de démonstration\\interprétation personnelle de l'œuvre}
\foldertitle{[Chapitre] III. Réseaux via découpages : notes manuscrites (08/06/1986) — lecture modernisée du dossier entier}

\begin{document}

\section*{Résumé}

\begin{resume}
Prenez une ficelle et des ciseaux. Un coup de ciseaux la coupe en deux
morceaux, deux coups en trois, $n$ coups en $n+1$. Supposons maintenant qu'on
ne vous montre pas la ficelle, mais seulement le résultat de toutes les coupes
possibles : quels points on peut couper ensemble, et, pour chaque coupe, quels
points tombent dans le même morceau. Pouvez-vous reconstituer la ficelle ? Ces
quarante feuillets, écrits les 8 et 9 juin 1986, répondent oui --- à
l'orientation près, puisque rien dans les coupes ne dit par quel bout
commencer --- et cherchent à faire de cette réponse une définition.

Le contexte est celui des deux chapitres précédents. Le premier (dossier 156-1)
décrivait une « forme » de dimension un --- un graphe dont les arêtes seraient
continues --- par ses lieux et ses segments ; le deuxième (dossier 156-2)
étudiait des espaces coupés par des « lieux », comme un cylindre l'est par des
cercles, et montrait que ces lieux s'ordonnent. Ce troisième chapitre propose
de tout reprendre à partir de la seule opération de \emph{découpage} : on se
donne quels lieux peuvent être coupés ensemble sans se gêner --- ils sont dits
en position « modérée » --- et, pour chaque ensemble fini de coupes, comment le
reste se partage en morceaux. L'espoir est que tout le reste en découle.

Le chapitre ne traite à fond que le cas le plus simple, celui de la ligne, mais
d'une ligne qui n'est pas forcément une droite : un ensemble ordonné où deux
points ne sont pas toujours comparables, comme les cercles d'un cylindre dont
deux peuvent se croiser. Il en dégage la notion de \emph{tronçon ordonné},
montre que l'ordre se retrouve à partir de la relation « $b$ est entre $a$ et
$c$ », puis celle-ci à partir des découpages. Vient ensuite l'étude des parties
« modérées » d'un tronçon : les réunions finies de morceaux et de points, ce
qu'une coupe en un nombre fini d'endroits peut produire. Elles forment une sorte
d'algèbre de Boole --- réunion, intersection, complémentaire --- à ceci près
qu'on ne peut combiner que des parties qui sont en position modérée l'une par
rapport à l'autre. L'image est celle de pièces de puzzle qu'on ne peut
assembler que si leurs bords se correspondent.

Le mot « modéré » n'est pas neutre : deux ans plus tôt, l'\emph{Esquisse d'un
programme} réclamait une « topologie modérée », débarrassée des pathologies de
la topologie générale. On en a aujourd'hui une version, la théorie des
structures o-minimales, dont le premier axiome dit justement que les parties
admissibles de la droite sont les réunions finies de points et d'intervalles.

La dernière page reprend la question de départ --- caractériser la ligne par
ses découpages --- et la laisse : « J'ai commencé à faire l'exercice, mais j'ai
l'impression que je n'ai pas la peine de l'expliciter jusqu'au bout. Ça doit
avoir été fait depuis longtemps (Hilbert ?) ». On voit là une manière de
travailler : poser une question, la pousser jusqu'à savoir qu'elle a une
réponse et de quelle nature, et ne pas s'attarder à la rédiger lorsqu'on la
croit déjà connue. Plusieurs questions restent ouvertes en chemin, posées
telles quelles --- « réciproque ? », « ??? » --- et c'est souvent là que le
lecteur trouvera à faire.

\keywords{betweenness relation, dense linear order, order topology, T0 space, interval algebra, tame topology, o-minimal structure, partial Boolean algebra}
\end{resume}

\section*{Le fil du dossier, et les conventions}

\begin{enumerate}
\item[1.] \emph{Les données d'un réseau par découpages} (pages 1 à 5), datées
du 8 juin 1986.
\item[2.] \emph{La ligne ordonnée et les tronçons} (pages 6 à 10).
\item[3.] \emph{Retrouver l'ordre à partir des découpages} (pages 10 à 16).
\item[4.] \emph{Parties modérées fermées} (pages 17 à 21).
\item[5.] \emph{La topologie d'un tronçon} (pages 21 à 24), datée du 9 juin.
\item[6.] \emph{Bord, tronçons et pseudo-tronçons} (pages 24 à 28).
\item[7.] \emph{Parties modérées et complémentaire modéré} (pages 29 à 31).
\item[8.] \emph{L'algèbre des parties modérées} (pages 31 à 33).
\item[9.] \emph{Relativisation et composantes connexes} (pages 34 à 39).
\item[10.] \emph{Retour au réseau} (page 40).
\end{enumerate}

Les feuillets portent sa pagination 1 à 40, qui coïncide avec celle des
archivistes. La page 1 porte le titre « Réseaux via découpages », la date du
8 juin 1986 et, dans un triangle, « GF III » : c'est le chapitre III de
\emph{Vers une géométrie des formes}, après « Vers une géométrie des formes
(topologiques) » (dossier 156-1, 5--7 juin) et « Réalisations topologiques des
réseaux » (dossier 156-2, 6--7 juin). Le lendemain du dernier feuillet, le
10 juin, s'ouvre au dossier 156-4 le chapitre IV, « Analysis situs (première
mouture) ».

\emph{Trois axiomatiques des réseaux.} Le dossier est la troisième tentative,
en quatre jours, de dire ce qu'est un réseau. Au dossier 156-1, un réseau est un
ensemble de lieux muni de segments et de leurs extrémités, soumis aux axiomes
F$_1$ à F$_6$ (pages 5 à 23 de ce dossier) ; ses lieux singuliers sont les
\emph{nœuds}, et les composantes des lieux réguliers ses \emph{lignes}. Au
dossier 156-2, le réseau est le but annoncé d'une étude des failles et des
mailles topologiques, et n'est pas atteint. Ici, les données premières sont les
paires de lieux en position modérée et les découpages. Les trois ne sont pas
fondues : chacune repart de zéro, et la seconde moitié de ce dossier abandonne
le réseau pour le seul tronçon. Deux liens sont pourtant explicites. Le
\emph{tronçon ordonné} de la page 8 est exactement ce que la proposition de la
page 9 du dossier 156-2 disait de l'ensemble des lieux d'une faille orientée,
moins l'existence d'un plus petit et d'un plus grand élément ; et la structure
de \emph{biordre} de la page 13 du dossier 156-1 --- une paire d'ordres
opposés --- est le « tronçon binordonné » de la page 11 ci-dessous. Le mot
« disjoints », enfin, a ici le sens qu'il avait au dossier 156-2 pour les
lieux d'une faille : distincts et comparables, et non d'intersection vide.

\emph{Conventions.} $\mathcal{L}$ désigne l'ensemble des lieux, puis, à partir
de la page 6, un tronçon ordonné\footnote{À partir de la page 21, la page écrit
$L$ en capitale droite. On garde $\mathcal{L}$ partout.}.
$\mathfrak{P}_{2}(\mathcal{L})$, $\mathfrak{P}_{f}(\mathcal{L})$ sont les parties
à deux éléments et les parties finies. $\mathcal{L}_{<a}$, $\mathcal{L}_{\leq a}$,
$]a,b[$, $[a,b]$, etc. ont leur sens ordinaire dans un ensemble ordonné, sans
qu'on suppose l'ordre total : $[a,b] = \{x \mid a \leq x \leq b\}$. Quatre
notions de bord coexistent dans les feuillets, et la page les note souvent du
même point suscrit ; on les distingue ainsi : $\dot{A}$, la frontière donnée
d'une composante dans l'axiomatique abstraite (page 3) ; $\partial T$, le bord
absolu d'un pseudo-tronçon, ensemble de ses points extrêmes (page 25) ;
$\mathrm{Fr}\, X$, la frontière modérée d'une partie modérée (pages 18 et 29) ;
$\mathrm{Fr}_{\mathrm{top}}\, T$, la frontière pour la topologie canonique
(page 25).

\emph{Ce que le dossier annonce et n'établit pas} : les axiomes des réseaux
par découpages eux-mêmes, au-delà des axiomes de finitude Dec$_1$, Dec$_2$ et de
l'axiome de la page 4 ; les conditions qui caractérisent la relation
d'intermédiarité d'un tronçon, et ses découpages (pages 16 et 40) ; la fin de
l'énoncé sur les parties localement fermées (page 20) ; les propriétés 1) à 8)
de l'algèbre des parties modérées (pages 32 et 33), énoncées sans
démonstration ; la sobriété de la topologie canonique (page 22).

\pagerange{1}{5}
\section*{Les données d'un réseau par découpages (pages 1 à 5)}

Le dossier s'ouvre sur une intention : « Je voudrais définir une autre
axiomatique des structures de réseau sur un ensemble $\mathcal{L}$ (ensemble
des lieux)\footnote{« autre axiomatique » est lu avec doute. Dans l'angle
supérieur gauche, deux ou trois mots biffés n'ont pas été lus.}. » Les données
sont au nombre de trois.

\begin{enumerate}
\item[(a)] Une partie $\mathfrak{S}_{0} \subset \mathfrak{P}_{2}(\mathcal{L})$,
dont les éléments sont les \emph{0-sphères} --- les paires de lieux
« modérées ». Si $\{a,b\} \in \mathfrak{S}_{0}$, on dit que $a$ et $b$ sont
\emph{disjoints}, ou en \emph{position relative modérée}. Par convention, un lieu
est en position modérée avec lui-même : la relation « $a = b$ ou
$\{a,b\} \in \mathfrak{S}_{0}$ » est réflexive et symétrique\footnote{La marge de
la page 1 le dit ; la page 2 écrit « $s \neq t$ ou $\{s,t\} \in
\mathfrak{S}_{0}$ », sous des notes marginales qui couvrent le début des lignes.
On lit $s = t$, comme dans la marge de la page 1 : avec $s \neq t$ la condition
serait toujours remplie.}. Une partie finie $T$ est \emph{modérée} si
$\mathfrak{P}_{2}(T) \subset \mathfrak{S}_{0}$, c'est-à-dire si ses éléments sont
deux à deux disjoints ; on note $\mathrm{Modf}(\mathcal{L})$ l'ensemble des
parties finies modérées.
\item[(b)] Pour tout $T \in \mathrm{Modf}(\mathcal{L})$, une relation
d'équivalence $R_{T}$ sur le \emph{complémentaire modéré}
\[
C(T) = \{t \in \mathcal{L} \smallsetminus T \mid T \cup \{t\} \text{ est
modérée}\},
\]
dont les classes sont les \emph{composantes connexes} de $C(T)$.
\item[(c)] Pour tout $T \in \mathrm{Modf}(\mathcal{L})$ et toute composante $A$
de $C(T)$, une partie $\dot{A} \subset T$, la \emph{frontière} de $A$.
\end{enumerate}

La page 1 dit ce que (a) devient dans l'axiomatique « par les segments » du
dossier 156-1 : deux lieux distincts sont disjoints si l'un d'eux est un nœud du
réseau, ou s'ils sont sur des lignes différentes, ou s'ils sont sur une même
ligne et forment le bord d'un segment\footnote{Le crochet de la page 1 n'est lu
qu'en partie : plusieurs mots sont surchargés ou illisibles, et « nœud »,
« lignes », « sur une même » sont des lectures douteuses. On n'en retient que
le sens d'ensemble, qui est celui des nœuds et des lignes définis à la page 23
du dossier 156-1.}. Sur une faille du dossier 156-2, c'est la relation entre
lieux de la page 8 de ce dossier, et un découpage par $n$ lieux deux à deux
disjoints y donne les $n+1$ failles recollées bout à bout de la page 9 : c'est
l'image à garder pour (b).

Il prévoit que les axiomes, encore à écrire, entraîneront que la donnée (c)
résulte de (b), et que (b) résulte de la seule donnée des $R_{\varepsilon}$ pour
$\varepsilon \in \mathfrak{S}_{0}$. La donnée essentielle serait alors la
relation de position modérée et, pour chaque 0-sphère $\varepsilon$, le
découpage de $C(\varepsilon)$ en composantes. C'est ce programme que les pages
10 à 15 réalisent pour un tronçon.

\emph{Parties modérées cofinies.} Les parties de la forme $C(T)$ sont dites
\emph{modérées cofinies}. Si une partie est à la fois finie modérée et modérée
cofinie, $U = S = C(T)$, alors $S \cup T$ est une partie finie modérée
maximale, puisque tout lieu modéré par rapport à $T$ est dans $S$. Il veut
qu'alors $S \cup T = \mathcal{L}$, que $\mathfrak{S}_{0} = \mathfrak{P}_{2}
(\mathcal{L})$, que les $R_{T}$ soient discrètes et les frontières vides. C'est
l'axiome suivant, accompagné de l'exemple qui le motive.

\begin{quote}
\textbf{Exemple 1} (page 4) : \emph{structure discrète.} Si $\mathfrak{S}_{0} =
\mathfrak{P}_{2}(\mathcal{L})$, alors $\mathrm{Modf}(\mathcal{L}) =
\mathfrak{P}_{f}(\mathcal{L})$ et $C(T) = \mathcal{L} \smallsetminus T$ ; on
prend $R_{T}$ discrète, les composantes sont des points et $\dot{A} =
\emptyset$.

\textbf{Axiome} (page 4). \emph{Si $T$ est une partie finie modérée maximale,
alors $T = \mathcal{L}$, et $\mathcal{L}$ est discret.}
\end{quote}

\emph{Finitude.} Viennent ensuite des hypothèses « draconiennes, peut-être
provisoires » :
\begin{enumerate}
\item[Dec$_{1}$] Toute partie modérée cofinie n'a qu'un nombre fini de
composantes connexes.
\item[Dec$_{2}$] Si $T \subset S$ sont finies modérées, de sorte que $U = C(S)
\subset V = C(T)$\footnote{La page écrit « $U \subset V$ i.e. $T \subset S$ ».
Seule l'implication $T \subset S \Rightarrow C(S) \subset C(T)$ est vraie en
général, et c'est la seule dont l'axiome se sert.}, toute composante $A$ de $U$
est contenue dans une composante $B$ de $V$, unique, et $\dot{A} \cap T \subset
\dot{B}$.
\end{enumerate}
Autrement dit, ajouter des coupes raffine le découpage, et une extrémité d'un
petit morceau qui était déjà une coupe reste une extrémité du grand morceau qui
le contient.

\pagerange{6}{10}
\section*{La ligne ordonnée et les tronçons (pages 6 à 10)}

\begin{quote}
\textbf{Exemple 2} (page 6). Soit $\mathcal{L}$ un ensemble ordonné filtrant
croissant et filtrant décroissant, sans plus grand ni plus petit élément,
localement filtrant croissant et décroissant, et divisible. On l'appelle une
\emph{ligne ordonnée}. On prend pour $\mathfrak{S}_{0}$ les paires totalement
ordonnées, c'est-à-dire les paires d'éléments distincts et comparables ; alors
$\mathrm{Modf}(\mathcal{L})$ est l'ensemble des parties finies totalement
ordonnées.
\end{quote}

Pour $T = \{t_{1} < \cdots < t_{n}\}$, un élément est modéré par rapport à $T$
s'il est comparable à chaque $t_{i}$ ; il se range alors dans la chaîne, et
\[
C(T) = \mathcal{L}_{<t_{1}} \sqcup \,]t_{1}, t_{2}[\, \sqcup \cdots \sqcup
\,]t_{n-1}, t_{n}[\, \sqcup \mathcal{L}_{>t_{n}} .
\]
On décrète que ces $n+1$ morceaux sont les composantes connexes ; ils sont non
vides, par divisibilité et faute d'extrémités, et leurs frontières sont
respectivement $\{t_{1}\}$, $\{t_{1}, t_{2}\}$, \dots, $\{t_{n-1}, t_{n}\}$,
$\{t_{n}\}$. La structure ne dépend donc que de l'ordre --- et même, puisque
rien ne change si on le renverse, de la paire formée de l'ordre et de son
opposé\footnote{La première ligne de la page 7 est lue en partie : « Cette
structure est donc déterminée par sa donnée binaire\dots{} elle ne dépend que du
binaire. » « Binaire » est son mot pour la paire d'ordres opposés --- le
« biordre » du dossier 156-1. Les mots « déterminée » et « donnée » sont
douteux.}. On notera que les morceaux ne sont pas totalement ordonnés :
$]t_{1}, t_{2}[$ contient en général des éléments incomparables entre eux.

\emph{Tronçons de la ligne} (page 7) : les parties $[s,t]$ ($s \leq t$),
$\mathcal{L}_{\leq s}$, $\mathcal{L}_{\geq s}$ et $\mathcal{L}$, de bords
$\{s,t\}$, $\{s\}$, $\{s\}$ et $\emptyset$. Une partie modérée serait une
réunion finie de tronçons disjoints deux à deux en position modérée, de bord la
réunion des bords\footnote{La condition « en position relative modérée » pour
deux tronçons n'est lue qu'en partie ; la page 26 la précisera : la réunion
des bords est totalement ordonnée. La marge de la page 7 porte une note
oblique lue par fragments.}. C'est l'ébauche de ce que les pages 17 à 21 et 29
à 31 reprendront.

La page 8 dégage alors la notion qui portera tout le dossier, en laissant
tomber l'hypothèse qu'il n'y a pas d'extrémités.

\begin{quote}
\textbf{Définition} (page 8). Un ensemble ordonné $\mathcal{L}$ est un
\emph{tronçon ordonné} si :
\end{quote}
\begin{enumerate}
\item[a)] pour tout $a$ tel que $\mathcal{L}_{>a} \neq \emptyset$,
$\mathcal{L}_{>a}$ est filtrant décroissant : si $x, y > a$, il existe $z$ avec
$x, y \geq z > a$ ;
\item[a')] dualement, si $x, y < a$, il existe $z$ avec $x, y \leq z < a$ ;
\item[b)] $\mathcal{L}$ est filtrant croissant et décroissant : pour $x, y$, il
existe $a, b$ avec $a \leq x, y \leq b$ ;
\item[c)] $\mathcal{L}$ est divisible : si $x < y$, il existe $z$ avec $x < z <
y$ ;
\item[d)] $\mathrm{card}\, \mathcal{L} > 1$.
\end{enumerate}

La condition d) est nécessaire, note la marge : a), b), c) sont satisfaites
par un ensemble à au plus un élément\footnote{Pour l'ensemble vide, b) n'est
satisfaite que si l'on ne demande pas qu'un ensemble filtrant soit non vide ;
c'est la convention implicite de la marge.}. Un tronçon est un \emph{segment}
s'il a un plus petit et un plus grand élément, un \emph{intervalle} s'il n'a ni
l'un ni l'autre, un \emph{demi-intervalle} s'il n'a que l'un des deux ; tout
tronçon est de l'une des trois sortes et d'une seule, et l'opposé d'un tronçon
est un tronçon de la même sorte, un demi-intervalle fermé à gauche devenant
fermé à droite. Les lieux d'une faille orientée du dossier 156-2 forment un
segment en ce sens : c'est la proposition de la page 9 de ce dossier. Une ligne
ordonnée est un intervalle.

On pose $\partial \mathcal{L}$ l'ensemble des éléments extrêmes (plus petit ou
plus grand), et
\[
\mathfrak{S}_{0}(\mathcal{L}) = \mathrm{Drap}_{2}(\mathcal{L}) = \{\{x,y\} \mid
x < y\},
\]
\[
\mathrm{Modf}(\mathcal{L}) = \mathrm{Drap}(\mathcal{L}) =
\{S \in \mathfrak{P}_{f}(\mathcal{L}) \mid S \text{ totalement ordonnée}\},
\]
les drapeaux de longueur deux et les drapeaux finis\footnote{« Drap » est son
abréviation de « drapeaux ». La page écrit $S \subset
\mathfrak{P}_{f}(\mathcal{L})$ pour $S \in \mathfrak{P}_{f}(\mathcal{L})$.}. Pour
$S = \{s_{1} < \cdots < s_{n}\}$, les composantes de $C(S)$ sont celles des
parties
\[
\mathcal{L}_{<s_{1}},\ ]s_{1}, s_{2}[,\ \dots,\ ]s_{n-1}, s_{n}[,\
\mathcal{L}_{>s_{n}}
\]
qui ne sont pas vides\footnote{La page met entre les $s_{i}$ des crochets qui se
lisent fermés, le premier étant surchargé ; les composantes sont dans $C(S)$,
qui ne contient pas les $s_{i}$ : ce sont les intervalles ouverts.} : seules
les deux extrêmes peuvent l'être, si $s_{1}$ est le plus petit élément ou
$s_{n}$ le plus grand. Pour $n = 0$ il y a une seule composante, $\mathcal{L}$ ;
pour $n = 1$, au plus deux. Les frontières sont les mêmes que plus haut.

La page 10 énonce alors le but : retrouver la structure binaire du tronçon, et
donc toute la structure (a), (b), (c), à partir de (a) seule et de (b) pour
$T \in \mathfrak{S}_{0}$. C'est, pour le tronçon, le programme de la page 3.

\pagerange{10}{16}
\section*{Retrouver l'ordre à partir des découpages (pages 10 à 16)}

\subsection*{L'ordre se lit sur l'intermédiarité}

Un \emph{tronçon binordonné} est un ensemble muni d'une paire d'ordres
opposés, chacun faisant de lui un tronçon ; c'est un tronçon dont on a oublié
l'orientation. Pour $a, b, c$ on note
\[
\mathcal{R}(a,b,c) \iff a < b < c \ \text{ou}\ a > b > c,
\]
« $b$ est strictement entre $a$ et $c$ », qui ne dépend pas du choix de l'ordre
dans la paire et qui est symétrique en $a$ et $c$. Sa version large,
$\overline{\mathcal{R}}(a,b,c) \iff a \leq b \leq c$ ou $a \geq b \geq c$,
s'en déduit : $\overline{\mathcal{R}}(a,b,c)$ équivaut à $\mathcal{R}(a,b,c)$
ou $b \in \{a,c\}$\footnote{La marge de la page 11 donne « ou $\{a,b,c\}$ est
de cardinal $\leq 2$ », précédé d'un mot illisible. Pris tel quel, c'est faux
pour $a = c \neq b$, où $\overline{\mathcal{R}}(a,b,a)$ est faux ; on écrit la
condition exacte.}. Dans la suite, « entre » sans adverbe est la relation
large.

\begin{quote}
\textbf{Lemme} (page 11). \emph{Une structure de tronçon binordonné sur
$\mathcal{L}$ est déterminée par la relation ternaire $\mathcal{R}$.}
\end{quote}

\emph{Démonstration.} D'abord $\mathcal{R}$ détermine
$\mathrm{Drap}_{2}(\mathcal{L})$ : deux éléments sont distincts et comparables
si et seulement si un troisième est strictement entre eux --- c'est la
divisibilité. Ensuite il existe une paire $\varepsilon = \{a,b\}$ d'éléments
comparables distincts : si $x \neq y$, b) donne $a \leq x, y \leq b$, et
$a \neq b$ puisque sinon $x = a = y$. On choisit l'ordre de la paire pour
lequel $a < b$ ; alors
\[
\mathcal{L}_{\geq b} = \{x \mid x = b \text{ ou } \mathcal{R}(a,b,x)\}, \quad
\mathcal{L}_{\leq a} = \{x \mid x = a \text{ ou } \mathcal{R}(b,a,x)\}, \quad
]a,b[ \,= \{x \mid \mathcal{R}(a,x,b)\}
\]
sont connus. Enfin $u \leq v$ si et seulement s'il existe $x \in
\mathcal{L}_{\leq a}$ et $y \in \mathcal{L}_{\geq b}$ avec $x \leq u \leq v
\leq y$ --- c'est la filtration b), qui fournit $x \leq u, a$ et $y \geq v, b$
---, et l'on a alors $x < y$. Il reste à exprimer $x \leq u \leq v \leq y$ par
l'intermédiarité, ce que fait le lemme suivant\footnote{La page 12 énumère les
conditions « $u$, $v$ sont entre $x$ et $y$ ; $u$ est entre $x$ et $v$ ; $v$
est entre $u$ et $y$ », réunies par un crochet. La page 13 les reprend en
lemme.}. \hfill$\square$

\begin{quote}
\textbf{Lemme} (page 13). \emph{Soient $x, y, u, v$ quatre éléments d'un ensemble
ordonné\footnote{La page écrit « trois éléments ».} tels que $x < y$, que $u$ et
$v$ soient entre $x$ et $y$, et que $u$ soit entre $x$ et $v$ (ou que $v$ soit
entre $u$ et $y$ : « une des deux suffit »). Alors $x \leq u \leq v \leq y$.}
\end{quote}

\emph{Démonstration.} Comme $x < y$, « $u$ entre $x$ et $y$ » donne $x \leq u
\leq y$, et de même pour $v$. Si $u$ est entre $x$ et $v$, ou bien $x \leq u
\leq v$, ou bien $x \geq u \geq v$ ; dans le second cas $u = x$ et $v \leq x
\leq v$, donc $u = v$ et la conclusion vaut encore\footnote{La page ramène
l'ensemble $\{x,y,u,v\}$ à une chaîne à deux, trois ou quatre éléments, après
un premier essai encadré et barré. L'argument est le même.}. \hfill$\square$

Le NB de la page 13 fait le compte : on a utilisé la divisibilité, et seulement
pour déduire $\mathrm{Drap}_{2}(\mathcal{L})$ de $\mathcal{R}$, et la filtration
b). Si l'on se donne à la fois $\mathrm{Drap}_{2}(\mathcal{L})$ et
$\mathcal{R}$, la filtration suffit pour retrouver l'ordre à l'opposé près
(page 14).

\subsection*{L'intermédiarité se lit sur les découpages}

Il reste à retrouver $\mathcal{R}$ à partir de $\mathrm{Drap}_{2}(\mathcal{L})$
et des découpages des $C(\varepsilon)$, $\varepsilon \in
\mathrm{Drap}_{2}(\mathcal{L})$\footnote{La page 14 écrit $C(\varepsilon) = \{t
\in \mathcal{L}_{\notin \varepsilon} \mid \{\varepsilon, t\} \in
\mathrm{Drap}_{3}(\mathcal{L})\}$, pour $t \in \mathcal{L} \smallsetminus
\varepsilon$ et $\varepsilon \cup \{t\}$.}. Pour $\varepsilon = \{a,b\}$ avec
$a < b$, un élément comparable à $a$ et à $b$ est plus petit que $a$, entre
eux ou plus grand que $b$, et
\[
C(\{a,b\}) = \mathcal{L}_{<a} \sqcup \,]a,b[\, \sqcup \mathcal{L}_{>b}.
\]

Pour que $\mathcal{R}(a,b,c)$ ait lieu, il faut d'abord que $J = \{a,b,c\}$ soit
un drapeau à trois éléments, ce que $\mathrm{Drap}_{2}(\mathcal{L})$ suffit à
dire. Pour chaque paire $\varepsilon \subset J$, soit $X_{\varepsilon}$ la
composante de $C(\varepsilon)$ qui contient le troisième élément de $J$. Si
$a < b < c$ :
\[
X_{\{a,b\}} = \mathcal{L}_{>b}, \qquad X_{\{b,c\}} = \mathcal{L}_{<b}, \qquad
X_{\{c,a\}} = \,]a,c[,
\]
d'où
\[
X_{\{a,b\}} \cap X_{\{b,c\}} = \emptyset, \qquad X_{\{a,b\}} \cap X_{\{c,a\}} =
\,]b,c[\, \neq \emptyset, \qquad X_{\{b,c\}} \cap X_{\{c,a\}} = \,]a,b[\, \neq
\emptyset,
\]
les deux dernières par divisibilité. L'élément du milieu est donc le seul
élément de $J$ pour lequel les deux paires qui le contiennent donnent des
composantes disjointes : $\mathcal{R}(a,b,c)$ se lit sur les découpages. « On
gagne ! » Et le programme de la page 3 est rempli pour le tronçon : la relation
de position modérée et les découpages des seules 0-sphères déterminent
l'ordre à l'opposé près, donc tout le reste.

Le nom actuel de $\mathcal{R}$ est celui de \emph{relation d'intermédiarité}
(\emph{betweenness})\footnote{Pour un ordre total, c'est la relation primitive
des axiomes d'ordre de M.~Pasch (1882) et de D.~Hilbert (\emph{Grundlagen der
Geometrie}, 1899, groupe II), axiomatisée pour elle-même par E.~V.~Huntington et
J.~R.~Kline (1917) ; les relations d'intermédiarité des ensembles partiellement
ordonnés ont été caractérisées par M.~Altwegg (1950). Les feuillets ne citent
personne ; la page 40 avance le nom de Hilbert, avec un point d'interrogation.
Que l'ordre soit déterminé à l'opposé près par l'intermédiarité est faux pour un
ensemble ordonné quelconque (deux composantes incomparables peuvent être
renversées indépendamment) ; c'est la filtration b) qui l'assure ici.}.

\subsection*{Remarques (page 16)}

(1) Soit $\mathcal{T}$ une sous-catégorie pleine de celle des ensembles
ordonnés telle que : a) tout objet est filtrant croissant et décroissant, de
cardinal $\neq 1$ (donc $\geq 2$) ; b) tout objet est divisible ; c) si
$\mathcal{L} \in \mathcal{T}$ et $\mathcal{L}_{>x} \neq \emptyset$, alors
$\mathcal{L}_{>x} \in \mathcal{T}$ ; c') l'énoncé dual. Il existe une plus
grande telle $\mathcal{T}$, stable par passage à l'ordre opposé, et ses objets
sont les tronçons ordonnés. C'est juste\footnote{La page ne démontre rien.
Une réunion de classes satisfaisant a)--c') les satisfait encore, d'où une plus
grande. Les tronçons forment une telle classe : si $\mathcal{L}$ est un tronçon
et $\mathcal{L}_{>x} \neq \emptyset$, $\mathcal{L}_{>x}$ est filtrant
décroissant par a), croissant par b), a au moins deux éléments par
divisibilité, et satisfait a') parce que $\mathcal{L}_{>x} \cap
\mathcal{L}_{<a} = \,]x,a[$ est filtrant croissant par a') dans $\mathcal{L}$.
Inversement un objet de la plus grande $\mathcal{T}$ satisfait b), c), d), et
a) parce que $\mathcal{L}_{>a}$ est dans $\mathcal{T}$, donc filtrant
décroissant. Ici « filtrant » entraîne « non vide ».} : c'est une
définition des tronçons par stabilité, et elle explique le choix des axiomes de
la page 8.

(2) Il faudrait expliciter quelles conditions $\mathcal{R}(a,b,c)$ doit
satisfaire pour provenir d'un tronçon binordonné, et de même pour la donnée de
$\mathfrak{S}_{0}$ et des $R_{\varepsilon}$. La question reste ouverte dans le
dossier ; la page 40 y revient.

\pagerange{17}{21}
\section*{Parties modérées fermées (pages 17 à 21)}

\emph{Sous-tronçons fermés.} Un \emph{sous-tronçon fermé} de $\mathcal{L}$ est
une partie de la forme $[a,b]$ avec $a < b$, $\mathcal{L}_{\leq a}$ avec $a$
non minimum, $\mathcal{L}_{\geq b}$ avec $b$ non maximum, ou $\mathcal{L}$ ;
c'est un tronçon pour l'ordre induit\footnote{Dans un ensemble filtrant
décroissant, un élément minimal est le plus petit ; « $a$ n'est pas le plus
petit élément » assure donc qu'il existe $x < a$, et $\mathcal{L}_{\leq a}$ a
au moins deux éléments.}. Un \emph{pseudo-tronçon fermé} est un sous-tronçon
fermé ou un point\footnote{La page 18 se sert de $\mathrm{Pstrf}(\mathcal{L})$,
« ensemble des ps-tronçons fermés », que la page 24 définira ; on prend cette
définition dès ici.}. Ses extrémités sont $a$ et $b$, $a$, $b$, aucune, ou le
point lui-même ; deux pseudo-tronçons fermés sont \emph{en position relative
modérée} si la réunion de leurs extrémités est totalement ordonnée\footnote{La
définition de la page 17 n'est lue qu'en partie ; on prend celle, explicite, de
la page 26. La page 17 note $\partial I$ l'ensemble des extrémités en ne
comptant pas celles qui sont des extrémités de $\mathcal{L}$ ($\partial
\mathcal{L}_{\leq a} = \{a\}$) ; mais $[a,b]$ avec $a$ minimum est aussi
$\mathcal{L}_{\leq b}$, et l'écriture n'est pas univoque. La convention qui
lève l'ambiguïté est celle de la page 29, rappelée plus bas ; elle ne change
rien à la position modérée, car un élément extrême de $\mathcal{L}$ est
comparable à tous.}. Une marge et quelques lignes biffées, lues par fragments,
évoquent déjà une topologie de base les $]a,b[$, $\mathcal{L}_{>a}$,
$\mathcal{L}_{<a}$ : c'est celle du 9 juin.

\begin{quote}
\textbf{Théorème} (page 18). \emph{Soit $X \subset \mathcal{L}$ une réunion
$X = \bigcup_{i \in I} T_{i}$ d'une famille finie de pseudo-tronçons fermés,
deux à deux disjoints et en position relative modérée. Alors la famille
$(T_{i})$ est déterminée par $X$ --- ce sont les composantes de $X$ --- et l'on
pose $\mathrm{Fr}\, X = \bigcup_{i} \mathrm{Fr}\, T_{i}$.} On dit que $X$ est une
\emph{partie modérée fermée}.
\end{quote}

L'unicité est juste\footnote{La page l'affirme sans démonstration. Deux
pseudo-tronçons fermés disjoints en position modérée sont comparables pour la
relation $T \ll T'$ de la page 27 (on le voit en rangeant leurs extrémités dans
une chaîne) ; on est donc dans la situation de la page 29, et l'on montre là
que les $T_{i}$ sont les pseudo-tronçons maximaux contenus dans $X$ : un
pseudo-tronçon contenu dans $X$ qui rencontrerait $T_{i}$ et $T_{j}$, $i \neq
j$, contiendrait, étant convexe, un point de l'interstice entre eux, qui n'est
pas dans $X$.}. La frontière modérée $\mathrm{Fr}\, X$ --- la page la note
$\dot{X}$ --- est vide si et seulement si $X = \emptyset$ ou $X =
\mathcal{L}$\footnote{Avec la convention de la page 29 : un élément extrême de
$\mathcal{L}$ n'est pas compté dans la frontière, sauf s'il forme à lui seul
une composante. La page écrit $\dot{X} = \bigcup \dot{T}_{i}$ ; elle note le
bord d'un point suscrit, après avoir écrit $\partial I$ à la page 17.}. Deux
parties modérées fermées $X$, $Y$ sont \emph{en position relative modérée} si
$\mathrm{Fr}\, X \cup \mathrm{Fr}\, Y$ est totalement ordonnée.

\begin{quote}
\textbf{Proposition} (page 19). \emph{Soient $X$, $Y$ deux parties modérées
fermées en position relative modérée. Alors $X \cap Y$ est une partie modérée
fermée, et}
\[
\mathrm{Fr}(X \cap Y) \subset \mathrm{Fr}\, X \cup \mathrm{Fr}\, Y .
\]
\emph{$X$ et $Y$ ont une borne supérieure $X \cup_{\mathrm{mod}} Y$ dans
l'ensemble des parties modérées fermées, et $\mathrm{Fr}(X \cup_{\mathrm{mod}}
Y) \subset \mathrm{Fr}\, X \cup \mathrm{Fr}\, Y$. Il en est de même d'une famille
finie $(X_{i})$ de parties modérées fermées deux à deux en position modérée.}
\end{quote}

La page n'en donne pas de démonstration ; l'énoncé est juste\footnote{Toutes
les extrémités sont dans une même chaîne. L'intersection de deux
pseudo-tronçons fermés dont les extrémités forment une chaîne est vide, ou le
pseudo-tronçon fermé qui va de la plus grande des deux extrémités inférieures à
la plus petite des deux supérieures. La borne supérieure s'obtient en
remplaçant chaque paquet de composantes de $X$ et de $Y$ qui se rencontrent de
proche en proche par le pseudo-tronçon fermé qui va de la plus petite de leurs
extrémités inférieures à la plus grande des supérieures ; l'argument de
maximalité donné pour le théorème de la page 18 montre que c'est la plus petite
partie modérée fermée contenant $X$ et $Y$. La page 19 écrit d'abord « borne
sup. \dots{} des parties modérées », puis, pour une famille, « dans l'ens. des
parties fermées modérées ». Deux marges, lues par fragments, portent un début
de démonstration et un exemple, peut-être $\{0, \frac{1}{2}, \frac{1}{3},
\dots\}$.}. La borne supérieure n'est pas la réunion :

\begin{quote}
\textbf{Exemple} (pages 19 et 20). $[a,b] \cup_{\mathrm{mod}} [b,c] = [a,c]$,
mais en général $[a,b] \cup [b,c] \neq [a,c]$. Si $a$ est le plus petit élément
et $c$ le plus grand, $[a,b] \cup [b,c]$ est l'ensemble des éléments
comparables à $b$, et $[a,c] = \mathcal{L}$. Ainsi $\mathcal{L}$ est totalement
ordonné si et seulement si, pour tout $b$,
\[
\mathcal{L}_{\leq b} \cup_{\mathrm{mod}} \mathcal{L}_{\geq b} \;(=
\mathcal{L}) \;=\; \mathcal{L}_{\leq b} \cup \mathcal{L}_{\geq b}.
\]
\end{quote}

\subsection*{Parties localement fermées}

Soient $X \supset Y$ deux parties modérées fermées en position relative modérée,
et
\[
C_{\mathrm{mod}}(Y,X) = \{x \in X \smallsetminus Y \mid x \text{ est en position
relative modérée avec } Y\},
\]
c'est-à-dire $\{x\} \cup \mathrm{Fr}\, Y$ totalement ordonné. Le théorème de la
page 20 affirme qu'il existe une plus petite partie modérée fermée $Z \subset X$
contenant $C_{\mathrm{mod}}(Y,X)$, que la marge note $X
\mathbin{-_{\mathrm{mod}}} Y$\footnote{La fin de la page 20 est surchargée
d'ajouts interlinéaires et n'est lue qu'en partie : on y devine une
caractérisation de $Z$ parmi les parties modérées fermées $T \subset X$
disjointes de $Y$ et en position modérée avec lui, et l'affirmation que $Z$ est
en position modérée avec $Y$. On ne la reconstitue pas. La page avait d'abord
écrit, puis barré, « la plus petite partie fermée modérée contenant $X
\smallsetminus Y$ » : $X \smallsetminus Y$ contient en général des éléments
incomparables aux extrémités de $Y$.}. La page 21 précise : si $F$ est la
frontière relative de $Y$ dans $X$, alors $F$ est une partie modérée fermée (une
chaîne finie), en position modérée avec $Y$, contenue dans $Z$, sans point
isolé de $Z$, et
\[
C_{\mathrm{mod}}(Y,X) = C_{\mathrm{mod}}(F, Z) .
\]
Sur un exemple totalement ordonné, $X = [0,3]$ et $Y = [1,2]$, on trouve
$C_{\mathrm{mod}}(Y,X) = [0,1[ \,\cup\, ]2,3]$, $Z = [0,1] \cup [2,3]$ et $F =
\{1,2\}$, et l'égalité est claire\footnote{La page n'en donne pas de
démonstration, et le sens exact de « frontière relative » n'est pas défini ;
l'exemple suggère $F = \mathrm{Fr}\, Y \smallsetminus \mathrm{Fr}\, X$. La lettre
$F$ est tracée en surcharge, et la mention « ne contient pas de pts isolés de
$Z$ » est lue avec doute.}. Une partie modérée localement fermée
serait ainsi de la forme $Z \smallsetminus F$ ; la définition n'est pas écrite.

\pagerange{21}{24}
\section*{La topologie d'un tronçon (pages 21 à 24)}

Le 9 juin, il munit un tronçon ordonné $\mathcal{L}$ d'une topologie, dite
\emph{canonique}, de base d'ouverts les $]a,b[$ ($a < b$), les
$\mathcal{L}_{<a}$ et les $\mathcal{L}_{>b}$. Ce sont bien les ouverts d'une
base : ce sont exactement les axiomes a) et a') des tronçons qui le
garantissent\footnote{La page ne le vérifie pas. Si $x \in \,]a,b[\, \cap
\,]c,d[$, a') donne $z$ avec $a, c \leq z < x$, la divisibilité $z'$ avec $z <
z' < x$, et de même au-dessus de $x$ : $x \in \,]z', w'[\, \subset \,]a,b[\,
\cap \,]c,d[$. Les autres cas sont semblables.}. Si $\mathcal{L}$ n'a pas de
plus petit élément, on peut omettre les $\mathcal{L}_{<a}$, réunions de
$]c,a[$, et dualement. Pour un tronçon totalement ordonné, c'est la topologie de
l'ordre.

Les points ne sont pas nécessairement fermés. Pour l'ensemble des lieux d'une
faille, la page relie l'adhérence d'un point à l'inclusion des lieux et en
déduit que $\overline{\{a\}} = \overline{\{b\}}$ entraîne $a = b$ : l'espace est
$T_{0}$\footnote{Les lignes de la fin de la page 21 sont très surchargées ;
seules les formules sont sûres, dont « $b \in \overline{\{a\}}$ i.e.
$\overline{\{b\}} \subset \overline{\{a\}}$ ». Le mot qui suit $T_{0}$ à la page
22, lu « prisotope », est douteux, comme le dérivé qui l'accompagne.}. Il ne
sait pas, dit-il, si c'est le cas pour tout tronçon, ni si la topologie canonique
est \emph{sobre} --- toute partie fermée irréductible aurait un point
générique\footnote{« sobre » est lu avec doute. Le terme est le sien, depuis SGA
4.}. La première question a une réponse négative : dans la droite réelle où
l'on dédouble l'origine en deux points incomparables $0$ et $0'$, placés l'un
et l'autre au-dessus des réels négatifs et au-dessous des réels positifs, on a
un tronçon ordonné, et tout ouvert de base qui contient l'un des deux points
contient l'autre ; l'espace n'est pas $T_{0}$\footnote{La vérification des
axiomes a) à d) est immédiate ; pour a') en $1$, par exemple, $\frac{1}{2}$
majore $0$ et $0'$ et est $< 1$. Le contre-exemple est le nôtre, et il ne dit
rien des lieux d'une faille, pour lesquels la page avance un argument qu'on n'a
pu lire.}. Lorsque $\mathcal{L}$ est totalement ordonné, en revanche, il est
séparé, et les $[a,b]$, $\mathcal{L}_{\geq a}$, $\mathcal{L}_{\leq a}$ sont
fermés ; la réciproque est posée en question.

Une seconde topologie, de base de fermés les $[a,b]$, $\mathcal{L}_{\geq a}$,
$\mathcal{L}_{\leq a}$, rendrait les « segments » fermés ; mais alors « on est
ennuyé, car $]a,b[$ n'a plus aucune raison d'être ouvert ! ». On peut prendre la
topologie engendrée par les deux familles d'ouverts --- $\mathcal{L}_{<a}$,
$\mathcal{L}_{>a}$, $]a,b[$ d'une part, $\complement \mathcal{L}_{\geq a}$,
$\complement \mathcal{L}_{\leq a}$, $\complement [a,b]$ d'autre part --- qui est
la moins fine des topologies plus fines que les deux\footnote{La page dit « la
moins fine des deux » ; il s'agit de leur borne supérieure. Elle écrit $[a,b]$
dans la seconde liste là où l'on attend $\complement [a,b]$, et la réunion des
deux listes n'est qu'une sous-base : il faut en prendre les intersections
finies.}. Les questions naturelles sont alors
\[
\overline{]a,b[} = [a,b] \ ? \qquad \overline{\mathcal{L}_{<a}} =
\mathcal{L}_{\leq a} \ ? \qquad \overline{\mathcal{L}_{>a}} = \mathcal{L}_{\geq a}
\ ?
\]
(pour $a$ non extrême dans les deux dernières). Le point crucial serait que tout
ouvert élémentaire contenant $a$ contienne un $]a, \varepsilon[$ ; c'est faux,
dit la page, déjà pour l'ensemble des lieux d'une faille : un croquis montre
$a \in \complement [x,y]$ avec $]a,b[ \,\subset [x,y]$\footnote{Le croquis de la
marge est une bande hachurée portant $x$, $a$, $b$, $y$. Pour des lieux non
totalement ordonnés, $a$ peut n'être pas entre $x$ et $y$ alors que tout lieu
strictement entre $a$ et $b$ l'est.}.

Il décide finalement (page 24) de garder la première topologie, et d'appeler
néanmoins, par abus de langage, \emph{sous-tronçons fermés} les $[a,b]$,
$\mathcal{L}_{\leq a}$, $\mathcal{L}_{\geq b}$ et $\mathcal{L}$, non réduits à un
point --- ceux de la page 17 --- bien qu'ils ne soient pas toujours fermés pour
elle.

\pagerange{24}{28}
\section*{Bord, tronçons et pseudo-tronçons (pages 24 à 28)}

\subsection*{Le bord absolu}

Pour un sous-tronçon fermé ou un point $T$, le \emph{bord absolu} est
\[
\partial T = \{t \in T \mid t \text{ est le plus petit ou le plus grand élément
de } T\},
\]
de cardinal $0$, $1$ ou $2$. Ainsi $\partial \{a\} = \{a\}$, $\partial [a,b] =
\{a,b\}$, et $\partial \mathcal{L}_{\geq a} = \{a\}$, $\partial
\mathcal{L}_{\leq a} = \{a\}$ si $\mathcal{L}$ n'a pas d'élément extrême du côté
ouvert\footnote{La page 25 écrit ces deux égalités « $\forall a \in L$ » ; si
$\mathcal{L}$ a un plus grand élément $\omega$, $\partial \mathcal{L}_{\geq a} =
\{a, \omega\}$. Elle remplace en revanche la valeur $\partial \mathcal{L} =
\emptyset$ de la page 24 par « $\partial L = \partial L$ (sic) » : le bord absolu
de $\mathcal{L}$ est $\partial \mathcal{L}$, vide seulement pour un intervalle.
La page 24, qui contenait ces définitions sous une première forme, est encadrée
et barrée.}. Le bord absolu se compare à la frontière topologique :
\[
\partial T = (\mathrm{Fr}_{\mathrm{top}}\, T \cap T) \cup (T \cap \partial
\mathcal{L}),
\]
réunion disjointe dès que $T$ n'est pas un point\footnote{La page encadre la
formule avec $\sqcup$. Pour $T = \{a\}$ avec $a \in \partial \mathcal{L}$, $a$
est dans les deux termes : il est dans la frontière, puisque $\mathcal{L}$ n'a
pas de point isolé, et dans $\partial \mathcal{L}$. Un premier énoncé encadré,
puis barré, traitait à part le cas $\mathrm{card}\, T = 1$. Démonstration : un
point de $T = [a,b]$ strictement entre $a$ et $b$ est intérieur, $]a,b[$ étant
ouvert ; $a$ est dans la frontière s'il n'est pas le plus petit élément de
$\mathcal{L}$, car tout ouvert de base qui le contient contient un point
$< a$, et il est intérieur s'il l'est, $\mathcal{L}_{<v}$ ($a < v < b$) étant
alors contenu dans $T$. Les autres cas sont semblables.} : un point de $T$ qui
est dans la frontière est une extrémité, et une extrémité qui n'est pas une
extrémité de $\mathcal{L}$ est dans la frontière.

\subsection*{Tronçons et pseudo-tronçons}

Un \emph{tronçon} de $\mathcal{L}$ est une partie de la forme $T' = T
\smallsetminus \alpha$, où $T$ est un sous-tronçon fermé et $\alpha \subset
\partial T$ :
\[
]a,b],\ [a,b[,\ ]a,b[ \ (a < b), \qquad \mathcal{L}_{<a} \ (a \neq
0_{\mathcal{L}}), \qquad \mathcal{L}_{>a} \ (a \neq \omega_{\mathcal{L}}),
\qquad \mathcal{L},
\]
avec les sous-tronçons fermés eux-mêmes\footnote{$0_{\mathcal{L}}$ et
$\omega_{\mathcal{L}}$ désignent le plus petit et le plus grand élément ; la
lettre lue $\omega$ est douteuse.}. Cette écriture est unique : $T$ est
l'adhérence de $T'$, au sens de la plus petite partie modérée fermée qui le
contient\footnote{La marge dit « c'est une adhérence, le plus petit fermé
contenant $T'$ ». Ce ne peut être l'adhérence topologique, dont la page 23 vient
de montrer qu'elle peut différer de $[a,b]$ ; c'est l'adhérence modérée de la
page 32, 6).}. Les tronçons $]a,b[$, $\mathcal{L}_{<a}$, $\mathcal{L}_{>a}$,
$\mathcal{L}$ sont dits \emph{ouverts} --- ce sont bien des ouverts de la
topologie canonique ---, et $]a,b]$, $[a,b[$ \emph{semi-ouverts}\footnote{La page
26 écrit « $[b,a[$ ».}. Attention : si $a$ est le plus petit élément, $[a,b[ =
\mathcal{L}_{<b}$ est ouvert ; $\mathcal{L}$ est à la fois ouvert et fermé.

Un \emph{pseudo-tronçon} est un tronçon ou un point. Pour $T' = T \smallsetminus
\alpha$ on pose $\partial T' = \partial T$, qui n'est plus, prévient la marge, le
bord absolu de $T'$ pour l'ordre induit\footnote{Une seconde formule, lue
« $\partial T' = \dot{T}' \sqcup \partial L$ », est douteuse ; on ne la reprend
pas. La marge ajoute que $\partial T' \neq \emptyset$ sauf si $T' = \mathcal{L}$,
avec une distinction de cas en partie illisible.}. Deux pseudo-tronçons $T$,
$T'$ sont \emph{en position relative modérée} si $\partial T \cup \partial T'$
est totalement ordonné.

\begin{quote}
\textbf{Proposition} (page 27). \emph{Si $T$, $T'$ sont des pseudo-tronçons en
position relative modérée, $T \cap T'$ est vide ou est un pseudo-tronçon, et
$\partial(T \cap T') \subset \partial T \cup \partial T'$.}
\end{quote}

« Immédiat » : les extrémités sont dans une chaîne, et l'intersection va de la
plus grande des extrémités inférieures à la plus petite des supérieures, une
extrémité étant prise si elle est dans les deux. Elle peut être un point
($]a,b] \cap [b,c] = \{b\}$) ou vide ($]a,b[ \,\cap\, [b,c]$).

\subsection*{L'ordre entre pseudo-tronçons}

Une relation d'ordre, qui n'est pas l'inclusion : pour deux pseudo-tronçons
disjoints et en position relative modérée, on écrit
\[
T \prec T' \iff t < t' \ \text{pour tous } t \in T,\ t' \in T' .
\]
Elle est transitive\footnote{La définition du corps de la page ne demande que
la condition sur les éléments ; la marge exige que $T$ et $T'$ soient en
position modérée et disjoints. Sans la position modérée, la suite est fausse :
dans le tronçon à origine dédoublée de la page 22, $]{-1}, 0[ \;\prec\; \{0'\}$
au sens du corps de la page, mais l'extrémité $0$ du premier et le point $0'$
sont incomparables.}. Si $T \prec T'$, $T$ est majoré, et a une borne
supérieure $a = \mathrm{Sup}\, T$, extrémité supérieure de son adhérence, avec
$T \subset \mathcal{L}_{\leq a}$ ; de même $b = \mathrm{Inf}\, T'$, et $a \leq
b$\footnote{La page écrit « $T$ a un plus grand élément $a$ » ; $]x,a[$ est
majoré et n'a pas de plus grand élément. $a$ et $b$ sont comparables par la
position modérée, et $b < a$ ferait se rencontrer $T$ et $T'$ près de $b$ ou
contredirait $T \prec T'$.}.

La position est \emph{de raccordement} si $a = b$ et si $a$ appartient à $T$ ou
à $T'$ --- à l'un des deux seulement, puisqu'ils sont disjoints : ainsi $]x,a]$
et $]a,y[$. On renforce $\prec$ en
\[
T \ll T' \iff T \prec T' \text{ sans raccordement} \iff \exists x \in
\mathcal{L},\ T \subset \mathcal{L}_{<x},\ T' \subset \mathcal{L}_{>x}.
\]
C'est la position relative \emph{propre}\footnote{Le premier indice est surchargé
à la page 27 et se lit $\leq x$ ; la page 28 écrit $<x$, qui est le bon.}.

\begin{quote}
\textbf{Proposition} (page 28). \emph{Soient $T \prec T'$, $a = \mathrm{Sup}\,
T$, $b = \mathrm{Inf}\, T'$. Sont équivalents :}
\end{quote}
\begin{enumerate}
\item[a)] \emph{$T$ et $T'$ ne se raccordent pas : si $a \in T$ ou $b \in T'$,
alors $a \neq b$ ;}
\item[b)] \emph{il existe $x$ avec $T \subset \mathcal{L}_{<x}$ et $T' \subset
\mathcal{L}_{>x}$ ;}
\item[c)] \emph{il existe un pseudo-tronçon $S$ avec $T \prec S \prec T'$, $T$
et $S$ se raccordant, ainsi que $S$ et $T'$.}
\end{enumerate}
\begin{quote}
\emph{De plus $S$ est unique, et $\partial S = \{a, b\}$ : c'est le
pseudo-tronçon d'extrémités $a$ et $b$, qui contient $a$ si et seulement si $a
\notin T$, et $b$ si et seulement si $b \notin T'$. On l'appelle le
\emph{pseudo-tronçon interstitiel} entre $T$ et $T'$.}
\end{quote}

\noindent Si $a = b$ n'appartient à aucun des deux, l'interstitiel est le point
$\{a\}$\footnote{La page décrit $S$ comme l'ensemble des $x$ tels que $T \subset
\mathcal{L}_{<x}$ et $T' \subset \mathcal{L}_{>x}$. Cela vaut si $a \in T$ et
$b \in T'$, ou si $\mathcal{L}$ est totalement ordonné, mais pas en général :
pour $T = \,]{-1}, 0[$ et $T' = \,]0, 1[$ dans le tronçon à origine dédoublée,
cet ensemble est $\{0, 0'\}$, qui n'est pas un pseudo-tronçon, alors que
l'interstitiel est $\{0\}$. Il faut lire « $x \geq a$ » lorsque $a \notin T$, et
« $x \leq b$ » lorsque $b \notin T'$. La marge note que b) entraîne $T \prec
T'$.}.

\begin{quote}
\textbf{Corollaire 1} (page 28). \emph{Soit $T$ un tronçon strictement minoré,
d'extrémité inférieure $c$. Il existe un unique pseudo-tronçon $S$ tel que $S
\prec T$, que $S$ et $T$ se raccordent, et que $S$ n'ait pas de minorant strict
: c'est $\mathcal{L}_{\leq c}$ si $c \notin T$, $\mathcal{L}_{<c}$ si $c \in
T$.}

\textbf{Corollaire 2.} \emph{On l'appelle le \emph{pseudo-tronçon interstitiel
gauche} de $T$ ; on définit de même, si $T$ est strictement majoré,
l'interstitiel droit.}
\end{quote}

\noindent La page décrit encore $S$ comme l'ensemble des minorants stricts de
$T$, avec la même réserve qu'à la proposition\footnote{« l'ens. des » est lu
avec doute. Pour $T = \,]0,1[$ dans le tronçon à origine dédoublée, les minorants
stricts sont $\mathcal{L}_{<0} \cup \{0, 0'\}$.}.

\pagerange{29}{31}
\section*{Parties modérées et complémentaire modéré (pages 29 à 31)}

Soit $\mathrm{Tr}(\mathcal{L})$ l'ensemble des pseudo-tronçons, ordonné par
$\ll$. Pour une partie finie totalement ordonnée $\Phi = \{T_{1} \ll \cdots \ll
T_{n}\}$ de $\mathrm{Tr}(\mathcal{L})$, on pose
\[
X_{\Phi} = \bigsqcup_{i} T_{i} ,
\]
réunion disjointe.

\begin{quote}
\textbf{Proposition} (page 29). \emph{L'application $\Phi \mapsto X_{\Phi}$, des
drapeaux finis de $(\mathrm{Tr}(\mathcal{L}), \ll)$ dans les parties de
$\mathcal{L}$, est injective.}
\end{quote}

En effet, les $T_{i}$ sont les pseudo-tronçons maximaux contenus dans
$X_{\Phi}$\footnote{La démonstration n'est pas sur la page. Un pseudo-tronçon est
convexe ; s'il était contenu dans $X_{\Phi}$ et rencontrait $T_{i}$ et $T_{j}$,
$i < j$, il contiendrait un point $x$ de l'interstitiel entre $T_{i}$ et
$T_{i+1}$, qui n'est dans aucun $T_{k}$. Sous le symbole $\mathrm{Drap}^{*}$, la
page porte « drapeaux de $\mathrm{Tr}(L)$, y inclus » suivi d'un mot
illisible ; on comprend que le drapeau vide est admis, et donne $X =
\emptyset$.}. Les parties de la forme $X_{\Phi}$ sont dites \emph{modérées} ; les
$T_{i}$ sont leurs \emph{composantes} --- « connexes ! », mais pas pour la
topologie canonique ---, et leur ensemble est noté $\pi_{0}(X)$. Les parties
modérées fermées des pages 17 à 21 sont celles dont toutes les composantes
sont fermées. La \emph{frontière} modérée est
\[
\mathrm{Fr}(X_{\Phi}) = \Bigl(\bigcup_{i} \partial T_{i}\Bigr) \smallsetminus
\varepsilon, \qquad \varepsilon = \{t \in \partial \mathcal{L} \cap X_{\Phi}
\mid \{t\} \notin \Phi\} :
\]
on retire de la réunion des bords, qui est un drapeau de $\mathcal{L}$, les
extrémités de $\mathcal{L}$ dont $X$ est un voisinage --- celles qui sont dans
$X$ sans y être une composante isolée. C'est une frontière relative à
$\mathcal{L}$, mais pas pour sa topologie, prévient la marge\footnote{La page
écrit $\mathcal{U}$ pour $X_{\Phi}$, indifféremment, aux pages 29 et 30.}.

Un $x \in \mathcal{L}$ est en position modérée avec $X$ s'il l'est avec chaque
composante, c'est-à-dire si $\{x\} \cup \mathrm{Fr}\, X$ est totalement ordonné
--- ce qui revient au même avec $\bigcup \partial T_{i}$, les éléments extrêmes
de $\mathcal{L}$ étant comparables à tous. Le \emph{complémentaire modéré}
\[
X' = C_{\mathrm{mod}}(X) = \{x \in \mathcal{L} \smallsetminus X \mid x \text{
est en position modérée avec } X\}
\]
est une partie modérée, et c'est la plus grande partie modérée $Y \subset
\mathcal{L} \smallsetminus X$ en position modérée avec $X$. Ses composantes
sont les $n - 1$ interstitiels entre $T_{i}$ et $T_{i+1}$, plus, s'ils existent,
l'interstitiel gauche de $T_{1}$ et l'interstitiel droit de $T_{n}$. Si $n'$ est
le nombre de composantes de $X'$, on a donc, pour $n \geq 1$\footnote{La page
écrit « minoré », « majoré » ; il faut lire « strictement » : $[0_{\mathcal{L}},
b]$ est minoré, mais n'a pas d'interstitiel gauche. Pour $n = 0$, $X = \emptyset$
et $X' = \mathcal{L}$, $n' = 1$. La marge de la page 30 étend la position
modérée à deux parties modérées : $\mathrm{Fr}\, X \cup \mathrm{Fr}\, Y$
totalement ordonné, ce qui est le cas si $X \subset Y$.} :
\[
n' = \begin{cases} n - 1 & T_{1} \text{ non strictement minoré, } T_{n} \text{
non strictement majoré,} \\ n & \text{un seul des deux,} \\ n + 1 & T_{1}
\text{ strictement minoré et } T_{n} \text{ strictement majoré,} \end{cases}
\]
d'où $|\mathrm{card}\, \pi_{0}(X) - \mathrm{card}\, \pi_{0}(X')| \leq 1$.

\pagerange{31}{33}
\section*{L'algèbre des parties modérées (pages 31 à 33)}

Sur $\mathcal{M} = \mathrm{Mod}(\mathcal{L})$, ensemble des parties modérées, il
y a trois structures : a) l'ordre d'inclusion ; b) la relation « $\{X, Y\}$
modérée » ; c) la partie $\mathcal{M}_{f} = \mathrm{Modfer}(\mathcal{L})$ des
parties modérées fermées. La marge demande si b) se déduit de l'ordre --- par
exemple, « $\{X,Y\}$ modérée ssi $X \wedge Y$ existe ??? » ---, et la question
reste sans réponse.

Les propriétés suivantes sont énoncées sans démonstration\footnote{On ne les
démontre pas non plus ici. Celles qu'on a vérifiées sur des exemples, et
qu'éclaire la description des composantes donnée plus haut, sont 1), 2), 5) et
la description de $X'$ ; les autres sont tenues pour les énoncés de la page.} :

\begin{enumerate}
\item[1)] $\mathcal{M}$ a un plus petit élément $\emptyset$ et un plus grand
$1 = \mathcal{L}$.
\item[2)] $X \subset Y \Rightarrow \{X, Y\}$ modérée.
\item[3)] Si $\{X,Y\}$ est modérée, $X \wedge Y$ et $X \vee Y$ existent dans
$\mathcal{M}$. Attention, note la marge : $X \vee Y$ n'est pas la réunion $X
\cup Y$.
\item[4)] Pour tout $X$, il y a un plus grand $Y$ parmi ceux tels que $\{X,Y\}$
soit modérée et $X \wedge Y = \emptyset$ : c'est $X'$, et c'est l'unique
élément tel que $\{X', X\}$ soit modérée, $X' \wedge X = \emptyset$ et $X' \vee
X = 1$. L'application $X \mapsto X'$ est une involution qui renverse l'ordre ;
$\{X',Y'\}$ est modérée si et seulement si $\{X,Y\}$ l'est, et alors $(X \wedge
Y)' = X' \vee Y'$, $(X \vee Y)' = X' \wedge Y'$. En fait, ajoute la marge,
$\{X,Y\}$ est modérée si et seulement si $\{X,Y'\}$ l'est.
\item[5)] $\emptyset, 1 \in \mathcal{M}_{f}$.
\item[6)] Toute $X$ est contenue dans une plus petite partie modérée fermée
$\overline{X}$, son \emph{adhérence} ; si $\{X,Y\}$ est modérée,
$\overline{X \vee Y} = \overline{X} \vee \overline{Y}$.
\item[7)] Si $\{X,Y\}$ est modérée, $\{\overline{X}, Y\}$ l'est aussi, donc
$\{\overline{X}, \overline{Y}\}$.
\item[8)] On définit l'intérieur $X^{\circ} = (\overline{X'})'$ et la frontière
$\mathrm{Fr}\, X = \overline{X} \wedge \overline{X'}$, de sorte que
$(\mathrm{Fr}\, X)' = X^{\circ} \vee X'^{\circ}$. Si $\{X,Y\}$ est modérée,
$\{\mathrm{Fr}\, X, \mathrm{Fr}\, Y\}$ l'est, et $\mathrm{Fr}\, X \cup
\mathrm{Fr}\, Y$ est modérée\footnote{La page écrit « $\mathrm{Fr}\, X \cup
\mathrm{Fr}\, X$ ». La marge voit là une propriété spécifique aux parties
modérées d'un objet de dimension un : la frontière y est une partie finie.}.
\end{enumerate}

Une marge de la page 32 dit l'essentiel : une famille de parties deux à deux
modérées engendre, pour $\wedge$, $\vee$ et $'$, une sous-algèbre de Boole,
finie si la famille l'est. En notation d'anneau,
\[
X + Y = (X \vee Y) \wedge (X \wedge Y)', \qquad X \cdot Y = X \wedge Y,
\]
\[
X \vee Y = X + Y + XY, \qquad X \leq Y \iff XY = X,
\]
avec $0 = \emptyset$ et $1 = \mathcal{L}$\footnote{La page écrit $X \vee Y = X + Y
- XY$, ce qui est la même chose en caractéristique deux.}.
Une structure de ce type --- un ensemble muni d'une relation réflexive et
symétrique de compatibilité, toute famille d'éléments deux à deux compatibles
étant contenue dans une algèbre de Boole --- s'appelle aujourd'hui une
\emph{algèbre de Boole partielle}\footnote{S.~Kochen et E.~Specker (1965--1967),
en logique quantique, où la compatibilité est la commensurabilité des
observables. Le rapprochement est le nôtre ; les axiomes 1) à 8) ne coïncident
pas avec les leurs. Lorsque $\mathcal{L}$ est totalement ordonné, toutes les
parties modérées sont deux à deux compatibles, et $\mathcal{M}$ est l'algèbre
de Boole des réunions finies de points et d'intervalles ouverts à extrémités
dans $\mathcal{L}$ --- une algèbre d'intervalles au sens de A.~Mostowski et
A.~Tarski (1939).}. Pour $\mathcal{L}$ totalement ordonné, c'est une algèbre de
Boole, celle des parties définissables de la droite dans un ordre dense : la
condition qui définit les structures \emph{o-minimales} est que les parties
définissables de la droite soient exactement de cette forme\footnote{L.~van den
Dries (1984), A.~Pillay et C.~Steinhorn (1986). Le livre de van den Dries,
\emph{Tame topology and o-minimal structures} (1998), prend son titre à la
« topologie modérée » de l'\emph{Esquisse d'un programme} (1984). Rien
n'indique que Grothendieck ait connu ces travaux en juin 1986 ; le rapprochement
est le nôtre.}.

Enfin, pour $X \in \mathcal{M}$, sont équivalents (page 33) : a) $X \subset
\mathrm{Fr}\, X$, c'est-à-dire $X^{\circ} = \emptyset$, ou encore $X'$ dense
($\overline{X'} = 1$) ; b) $X = \mathrm{Fr}\, X$, c'est-à-dire $X$ fermé et
rare, $X'$ ouvert dense ; c) $X$ est une partie finie de $\mathcal{L}$. Les
parties rares sont les chaînes finies\footnote{Un premier énoncé de cette
équivalence, sous le titre « $X$ est rare si $X^{\circ} = \emptyset$ », est
barré de quatre traits et repris aussitôt. La page note la frontière d'un point
suscrit en a) et b).}. Une phrase en partie illisible demande quelles
conditions imposer à un ensemble ordonné $\mathcal{M}$ satisfaisant 1) à 7)
pour qu'il soit de cette forme.

\pagerange{34}{39}
\section*{Relativisation et composantes connexes (pages 34 à 39)}

\subsection*{Relativiser à une partie}

Pour $A \in \mathcal{M}$, soit $\mathcal{M}_{A} = \mathcal{M}_{\leq A} = \{Y
\in \mathcal{M} \mid Y \subset A\}$, avec l'ordre induit et la relation
« modérée » induite ; $X \in \mathcal{M}_{A}$ est \emph{fermé dans $A$} si $X =
\overline{X} \wedge A$, c'est-à-dire s'il existe $Z \in \mathcal{M}_{f}$ avec
$\{Z, A\}$ modérée et $X = Z \wedge A$. Les conditions 1) à 7) restent
satisfaites, et la page les vérifie une à une, en admettant celles de
$\mathcal{M}$ :
\begin{enumerate}
\item[1)] $0_{\mathcal{M}_{A}} = \emptyset$ et $1_{\mathcal{M}_{A}} =
A$\footnote{La page écrit « $1_{\mathcal{M}_{A}} = X$ ».} ;
\item[2)] est triviale ;
\item[3)] les bornes $X \wedge Y$, $X \vee Y$ calculées dans $\mathcal{M}$ sont
majorées par $A$, et sont les bornes dans $\mathcal{M}_{A}$ ;
\item[4)] le complémentaire de $X$ dans $\mathcal{M}_{A}$ est $X' \wedge A$ :
un $Y \subset A$ modéré avec $X$ et tel que $X \wedge Y = \emptyset$ est majoré
par $X'$ et par $A$\footnote{La page écrit « $Y \in \mathcal{M}_{Y}$ » pour
$\mathcal{M}_{A}$.}. Plus conceptuellement (page 35), les opérations de
$\mathcal{M}_{A}$ sur une famille de parties deux à deux modérées, et modérées
avec $A$, sont celles de $\mathcal{M}$ suivies du produit par $A$ : on travaille
dans l'idéal principal engendré par $A$ dans une algèbre de Boole, qui est une
algèbre de Boole d'unité $A$. Et si $\{X,Y\}$ est modérée, $\{A \wedge X', A
\wedge Y'\}$ l'est aussi ;
\item[5)] est triviale ;
\item[6)] l'adhérence de $X$ dans $\mathcal{M}_{A}$ est $\overline{X} \wedge A$,
et pour $\{X,Y\}$ modérée,
\[
\overline{X \vee Y}^{A} = \overline{X \vee Y} \wedge A = (\overline{X} \vee
\overline{Y}) \wedge A = (\overline{X} \wedge A) \vee (\overline{Y} \wedge A) =
\overline{X}^{A} \vee \overline{Y}^{A} ;
\]
\item[7)] si $\{X,Y\}$ est modérée, $\{X, Y, A\}$ l'est, donc $\{\overline{X},
Y, A\}$, donc $\{\overline{X}^{A}, Y\}$. « OK ».
\end{enumerate}

\subsection*{Connexité}

On dit que $\mathcal{M}$ est \emph{connexe} si les seuls éléments à la fois
ouverts et fermés sont $\emptyset$ et $1$. Lorsque les éléments ouverts et
fermés engendrent une algèbre de Boole finie, les \emph{composantes connexes}
sont ses atomes, les éléments non nuls minimaux parmi les ouverts
fermés\footnote{Le NB de la page 36 est surchargé et n'est lu qu'en partie ; on
y devine que les ouverts fermés sont deux à deux modérés, ce qui est nécessaire
pour qu'ils engendrent une algèbre de Boole. La page suppose $0 \neq 1$.}.

\begin{quote}
\textbf{Exemple} (pages 37 et 38). \emph{Soit $A = X_{\Phi} = \bigsqcup T_{i}$
une partie modérée de $\mathcal{L}$. Les composantes connexes de
$\mathcal{M}_{A}$ sont les pseudo-tronçons $T_{i}$.}
\end{quote}

\emph{Démonstration.} On a $A = \bigvee T_{i}$ dans $\mathcal{M}$, et $T_{i}
\wedge T_{j} = \emptyset$ pour $i \neq j$. Il suffit de voir que chaque $T_{i}$
est ouvert et fermé dans $A$, et connexe. a) $T_{i}$ est fermé dans $A$ :
$\overline{T_{i}} \cap A = T_{i}$, parce que l'adhérence modérée d'un
pseudo-tronçon est le pseudo-tronçon fermé associé, dont les extrémités
ajoutées ne sont pas dans $A$ ; les autres $T_{j}$ l'étant aussi, $T_{i}$ est
ouvert. b) $T$ est connexe : on regarde $T$ comme un tronçon en lui-même, ce qui
ramène au cas $T = \mathcal{L}$ (puisque $(\mathcal{M}_{\mathcal{L}})_{T} =
\mathcal{M}_{T}$). Une partie $X$ ouverte et fermée a une frontière
$\overline{X} \wedge \overline{X'} = X \wedge X' = \emptyset$ ; la frontière
modérée de $X = \bigsqcup T_{j}$ est vide si et seulement si chaque $T_{j}$ est
$\mathcal{L}$ ; donc $X = \emptyset$ ou $X = \mathcal{L}$\footnote{La page écrit
que la frontière est la réunion des $\partial T_{j}$ et que $\partial T_{j} =
\emptyset$ entraîne $T_{j} = \mathcal{L}$. Pour le bord absolu, c'est faux si
$\mathcal{L}$ est un segment ($\partial \mathcal{L} \neq \emptyset$) ; c'est
vrai pour la frontière modérée de la page 29, qui ôte les extrémités de
$\mathcal{L}$. On identifie la frontière abstraite de la page 33 à celle-là,
comme le fait la page. Le NB de la page 38 sur l'adhérence n'est lu qu'en
partie ; sous l'accolade de a), sa main ajoute « modérés bien sûr ! ».}.
\hfill$\square$

\subsection*{Subdivision}

\begin{quote}
\textbf{NB} (page 39). \emph{Soit $(X_{\alpha})$ une famille finie de parties
modérées deux à deux modérées, de sorte que $S = \bigcup_{\alpha} \mathrm{Fr}\,
X_{\alpha}$ est une chaîne $t_{1} < \cdots < t_{n}$. L'algèbre de Boole
engendrée par les $X_{\alpha}$ a pour atomes, parmi}
\[
\mathcal{L}_{<t_{1}}\ (\text{si } t_{1} \notin \partial \mathcal{L}),\
\{t_{1}\},\ ]t_{1}, t_{2}[,\ \{t_{2}\},\ ]t_{2}, t_{3}[,\ \dots,\ \{t_{n}\},\
\mathcal{L}_{>t_{n}}\ (\text{si } t_{n} \notin \partial \mathcal{L}),
\]
\emph{ceux qui sont nécessaires pour engendrer les $X_{\alpha}$ ; elle est
donc connue dès qu'on connaît $S$.}
\end{quote}

C'est la \emph{subdivision} du tronçon par la chaîne $S$\footnote{La page écrit
$]t_{1}, t_{3}[$ pour $]t_{2}, t_{3}[$. Les lignes qui précèdent la liste sont
barrées et lues par fragments ; on y lit que les composants sont des
pseudo-tronçons. Que l'algèbre engendrée soit contenue dans celle des réunions
de ces $2n + 1$ morceaux est ce que la page affirme ; elle peut en être une
sous-algèbre stricte, d'où la formulation.}. Si $\mathcal{L}$ n'est pas
totalement ordonné, ces morceaux ne recouvrent pas $\mathcal{L}$ comme ensemble
: un élément incomparable à $t_{1}$ n'est dans aucun. Leur borne supérieure dans
$\mathcal{M}$ est pourtant $1$ --- c'est l'avertissement de la page 32,
$X \vee Y \neq X \cup Y$.

\pagerange{40}{40}
\section*{Retour au réseau (page 40)}

La dernière page revient à la question de départ, sous le titre
« Caractérisation d'un intervalle ordonné, en tant que réseau » : caractériser
un intervalle ordonné par la donnée de
\[
S \subset \mathfrak{P}_{2}(\mathcal{L}), \quad \text{qui doit être }
\mathrm{Drap}_{2}(\mathcal{L}),
\]
et, pour $\varepsilon = \{a,b\} \in S$, de la décomposition de
\[
\complement_{\mathrm{mod}}\, \varepsilon = \{x \in \mathcal{L} \mid \{a,x\},
\{b,x\} \in S\}
\]
en composantes connexes, par une relation $R_{\varepsilon}$. C'est exactement
la donnée (a) des 0-sphères et la donnée (b) restreinte aux 0-sphères de la
page 3 ; les pages 10 à 15 ont montré qu'elle détermine l'intervalle, et la
remarque (2) de la page 16 demandait quels axiomes elle doit satisfaire.

« J'ai commencé à faire l'exercice, mais j'ai l'impression que je n'ai pas la
peine de l'expliciter jusqu'au bout. Ça doit avoir été fait depuis longtemps
(Hilbert ?) dans le cas totalement ordonné (et divisible) (en termes de la
relation : $s$ est entre $a$ et $b$\dots) ». Dans le cas totalement ordonné, la
caractérisation par l'intermédiarité est en effet classique ; dans le cas
partiellement ordonné, qui est celui des lieux d'une faille et donc celui qui
l'intéressait, elle l'est moins\footnote{Voir la note sur la relation
d'intermédiarité, pages 10 à 16 : Huntington et Kline (1917) pour un ordre
total, Altwegg (1950) pour un ordre partiel. Ni les uns ni les autres ne
traitent la donnée des découpages $R_{\varepsilon}$, qui est celle de la page.
Le titre corrige un premier « réseau intervalle », biffé.}. Le dossier s'arrête
là. Le lendemain, le chapitre IV recommence autrement, par l'« Analysis situs ».

\end{document}
